\bookchapter{Chapter 35: Children and Phones — An Educational Guide for the Digital Age}

\chapterintro{This chapter doesn't paint phones as “monsters,” nor does it advocate “the earlier the smarter.” We'll go step by step by age: giving devices, setting rules, teaching literacy, and protecting health.}

\booksection{I. At What Age Should Kids Get a Phone?}

There's no standard answer, but there are age baselines. Giving a phone too early — the child's prefrontal cortex isn't fully developed, so self-control can't keep up with algorithms' “temptation design.” Too late — the child may become “socially disconnected” from peers. Here are age-based recommendations:

\begin{itemize}
  \item \textbf{Ages 3-6}: Try to avoid phones entirely. If needed, no more than 15 minutes per session, with parent-selected educational cartoons, and must be\textbf{co-viewed with a parent}. Screen time has the greatest impact on vision and attention development at this stage.
  \item \textbf{Ages 7-9}: Can "borrow" a parent's device for online classes, audiobooks, and research. Start establishing "request before use, return after" rules. No dedicated phone.
  \item \textbf{Ages 10-12}: Consider a "feature phone" or kids' smartwatch for school commute communication. If a smartphone is truly needed, enable strict parental controls.
  \item \textbf{Ages 13+}: Entering middle school, social needs increase. A smartphone can be given, but with a "usage agreement" and regular reviews.
\end{itemize}

\begin{tcbinfo}{The "Good Enough" Principle}
A child's first phone doesn't need to be a flagship. “Excess performance” processors often end up running games. Choose a mid-range phone with adequate battery, eye-friendly screen, or an official refurbished model — easier to manage. Budget reference: 1,500-2,500 yuan covers most students' needs.

\end{tcbinfo}

\booksubsection{Smartwatch as a Transition}

If you just want to solve the “can I reach them after school” problem, a kids' smartwatch is better than a phone. It can only make calls, send voice messages, and track location — no browser or short videos, naturally limiting “addiction gateways”:

\begin{itemize}
  \item \textbf{XiaoTianCai (Little Genius) Phone Watch}: Most complete ecosystem, peer social "tap to add friend" is its core selling point, but the proprietary app ecosystem is relatively closed.
  \item \textbf{Huawei Kids' Watch}: Good integration with Huawei phones, accurate positioning, supports video calls, high value for money.
  \item \textbf{Xiaomi Kids' Watch}: Affordable, complete basic features, suitable for budget-conscious families.
\end{itemize}

When buying a watch, note: choose one that supports\textbf{4G full network compatibility}models (2G/3G are being phased out), and check carrier kids' plan rates.

\booksection{II. Complete Parental Control Guide}

Before giving the device to the child, the first thing isn't turning it on — it's setting up parental controls. Both platforms have native solutions:

\booksubsection{iOS: Screen Time}

\begin{itemize}
  \item \textbf{App Limits}: Settings > Screen Time > App Limits — set daily time limits per app category (e.g., games 30 minutes).
  \item \textbf{Downtime}: Set a daily "lights out" period (e.g., 22:00-7:00) — all apps gray out, only phone calls available.
  \item \textbf{Content Restrictions}: Filter movies, web content, and app installs by age, and disable in-app purchases.
  \item \textbf{Family Sharing}: Use the parent's iPhone to remotely manage the child's device, approve download requests, and view usage reports.
\end{itemize}

\booksubsection{Android: Google Family Link / Brand Parental Controls}

\begin{itemize}
  \item \textbf{Google Family Link}: Google's official solution — remotely manage device usage time, lock devices, view activity reports. Availability on domestic Android varies by model.
  \item \textbf{Xiaomi Kids Space}: Settings > Kids Space — set usage time, whitelist apps, requires parent password to exit.
  \item \textbf{Huawei Student Mode}: Settings > Healthy Phone Usage > Student Mode — restrict usage times and apps, supports remote management.
  \item \textbf{OPPO/vivo Kids Mode}: Enable in "Settings" or the dedicated "Kids Space" app — similar principle.
\end{itemize}

\begin{center}
\begin{tabularx}{\linewidth}{|X|X|X|}
\hline
\textbf{Feature} & \textbf{iOS Screen Time} & \textbf{Android Family Link / Brand Solutions} \tabularnewline
\hline
App Limits & Supported, by category & Supported, per app \tabularnewline
Downtime & Supported & Supported (brand solutions vary slightly) \tabularnewline
Content Filtering & System-level, more thorough & Depends on brand, varies \tabularnewline
Remote Management & Family Sharing & Family Link / Brand Cloud \tabularnewline
App Install Approval & Requires parent approval & Family Link supports it \tabularnewline
\hline
\end{tabularx}
\end{center}

\begin{tcbwarn}{Parental Controls Aren't "Once and Done"}
Kids can always find ways to bypass them (search “how to crack Screen Time” online — tons of tutorials). Technical measures are just the baseline — the real defense is parent-child communication and trust. I recommend reviewing usage reports together with your child monthly, rather than unilaterally monitoring.

\end{tcbwarn}

\booksection{III. Youth Mode}

Major content platforms have built-in “Youth Mode” — when enabled, it limits usage time, filters inappropriate content, and disables in-app purchases and tipping. This is a legally mandated feature, but\textbf{parents must actively enable it and set a password}:

\begin{itemize}
  \item \textbf{Douyin/Kuaishou Youth Mode}: 40 minutes daily (disabled 22:00-6:00), only curated science and educational videos shown — no livestreaming, no tipping, no recharging.
  \item \textbf{WeChat Youth Mode}: Limits WeChat Pay (requires parent-authorized amount), restricts Official Account and Channels content, Shake and Nearby features are disabled by default.
  \item \textbf{Game Anti-Addiction System}: National-level real-name authentication + time limits. Minors can only log into games on Friday, Saturday, Sunday, and public holidays from 20:00-21:00, for 1 hour per day.
\end{itemize}

\textbf{Don't Use Parents' IDs for Real-Name Authentication to “Borrow Accounts”}\par

Many kids use parents' IDs to register game accounts and bypass anti-addiction. This violates user agreements and removes protections. Parents should safeguard their payment passwords and identity information — don't “lend accounts” to children for convenience.

\booksection{IV. Digital Literacy Education}

Technology can block 90\% of risks — the remaining 10\% depends on the child's judgment. Rather than “no internet,” it's more important to teach them “how to use the internet”:

\begin{itemize}
  \item \textbf{Privacy protection}: Don't share real name, school, or home address with strangers. Don't send personal photos, ID cards, or family group photos. Keep location services off by default, only enable when needed.
  \item \textbf{Cyberbullying Recognition and Response}: Being insulted, ostracized, or having photos maliciously edited are all forms of bullying. Tell your child: don't retaliate, don't delete evidence, tell a parent immediately. On WeChat, long-press the message to report — the platform will handle the bully's account.
  \item \textbf{Content Criticism}: Teach children to distinguish "facts" from "opinions," recognize clickbait and sensationalism. Understand that algorithms only push what you want to see, not the whole picture. Treat "too good to be true" ads as always fake.
  \item \textbf{Cybersecurity}: Don't click unfamiliar links, don't scan QR codes from unknown sources, don't download apps from non-official stores. Don't share passwords with anyone (including "best friends"). Anyone asking for verification codes is a scammer.
\end{itemize}

\begin{tcbpractice}{Hands-On: Sign a "Phone Usage Agreement" with Your Child}
Writing rules down, both parties signing, and posting them prominently is far more effective than verbal agreements. Recommended contents:

\begin{itemize}
  \item Usage times: e.g., "Available 19:30-20:30 after homework is done."
  \item No-phone zones: at the dinner table, in the bedroom before sleep, in class — phone put away.
  \item Content boundaries: don't install apps without telling parents, don't private chat with strangers.
  \item Violation consequences: e.g., "One violation halves next day's usage time."
  \item Parent commitment: don't "scroll" on the phone in front of the child beyond agreed time — lead by example.
\end{itemize}

\end{tcbpractice}

\booksection{V. Healthy Usage Habits}

Children's eyes, neck, and sleep are still developing — phone damage is greater than for adults:

\begin{itemize}
  \item \textbf{Screen Time Recommendations}: Ages 3-6: no more than 1 hour daily; ages 6-12: no limit for educational content, no more than 1 hour for entertainment; ages 12+: negotiate with the child, but entertainment no more than 2 hours.
  \item \textbf{Eye Protection Habits}: Follow the "20-20-20 rule" — every 20 minutes of screen time, look at something 20 feet (about 6 meters) away for 20 seconds. No phone 1 hour before bed — blue light suppresses melatonin and affects sleep.
  \item \textbf{Posture Management}: Keep the phone screen at eye level, don't hunch over. Get up and move after every 30 minutes of continuous use.
  \item \textbf{Family Shared Rules}: Set "no-phone zones" at the dining table, bedroom, and study. During family movie time and outdoor activities, everyone puts their phone away — making "putting down the phone" a family culture, not a punishment for the child.
\end{itemize}

\begin{tcbinfo}{Parents Are the Best "Firewall"}
Research shows a strong positive correlation between parents' phone usage time and children's. Demanding that children stay off their phones while scrolling short videos until late at night yourself makes the rules ineffective. To change your child, start by putting down your own phone.

\end{tcbinfo}

\booksection{VI. Digital Legacy and Children}

This section is easily overlooked but very important. Children's digital assets (accounts, photos, stored value, game items) also face inheritance and privacy issues:

\begin{itemize}
  \item \textbf{Inheritance of Children's Digital Assets}: When a child unfortunately passes away, the inheritance process for social accounts, game accounts, and cloud photos is the same as for adults — requires death certificate and proof of kinship. Parents should set up "legacy contacts" or "emergency contacts" for children's important accounts in advance (e.g., Apple Legacy Contact, WeChat account inheritance).
  \item \textbf{Digital Privacy of Minors}: The Personal Information Protection Law stipulates that processing personal information of minors under 14 requires guardian consent, with dedicated privacy rules. Parents should regularly review the privacy policies of apps their children use frequently, and turn off "personalized recommendations" and "data sharing" toggles.
  \item \textbf{Account Registration Recommendations}: Accounts registered for children should use parent-controlled email/phone numbers for unified management during the minor stage; account transfer can be done when the child reaches adulthood.
\end{itemize}

\begin{tcbwarn}{The Other Side of "Children's Digital Legacy"}
Children are being “recorded” from birth — growth photos posted by parents, videos from kindergarten, health data from smartwatches. Once uploaded, this data is hard to completely delete. Before posting, think: Will the child feel embarrassed by this photo in ten years? Respecting your child's digital identity starts with restraining from “sharenting.”

\end{tcbwarn}

\begin{tcbsummary}{Chapter Summary}
\begin{itemize}
  \item There's no standard age to give a phone, but ages 3-6 should avoid it, and before 13 use a kids' watch or feature phone as transition
  \item "Good enough": choose a budget phone, not a flagship — excess performance equals excess temptation
  \item iOS Screen Time and Android Family Link/brand parental controls are two technical defense lines — use them together for better protection
  \item Youth Mode and game anti-addiction are legally required, but parents must actively enable them and guard the passwords
  \item Digital literacy is more important than technical blocking: teach kids to recognize bullying, protect privacy, and critically evaluate content
  \item Healthy habits start with parents — when you put down your phone, your child will too
  \item Children's digital legacy and privacy are emerging issues — set up "legacy contacts" and restrain from sharenting
\end{itemize}

\end{tcbsummary}
